\documentclass[a4paper,11pt]{article}

\usepackage[top=2.5cm, bottom=2.5cm, left=2.5cm, right=2.5cm]{geometry}
\usepackage{amsmath, amssymb}
\usepackage{fancyhdr}
\usepackage{titlesec}
\usepackage{times}

% ── Header / Footer ──
\pagestyle{fancy}
\fancyhf{}
\renewcommand{\headrulewidth}{0pt}
\lhead{\textsf{Release~19}}
\chead{\textsf{\textbf{37}}}
\rhead{\textsf{3GPP TR 25.996 V19.0.0 (2025-09)}}

\cfoot{\textbf{\textsf{3GPP}}}

% ── Section formatting ──
\titleformat{\section}[block]{\Large\bfseries}{}{0pt}{}[\vspace{-0.3em}\rule{\textwidth}{0.4pt}]

% ── Equation numbering style ──
\renewcommand{\theequation}{A-\arabic{equation}}

\begin{document}

\section*{Annex A:\\Calculation of circular angle spread}
\addcontentsline{toc}{section}{Annex A: Calculation of circular angle spread}

\noindent\rule{\textwidth}{0.4pt}

\vspace{1em}

For a signal with \textit{N} multi-paths, each has \textit{M} sub-paths, the conventional angle spread calculation for the composite signal is given by

\begin{equation}
\sigma_{AS} = \sqrt{\frac{\displaystyle\sum_{n=1}^{N}\sum_{m=1}^{M}\bigl(\theta_{n,m,\mu}\bigr)^{2}\cdot P_{n,m}}{\displaystyle\sum_{n=1}^{N}\sum_{m=1}^{M}P_{n,m}}}
\label{eq:A1}
\end{equation}

where $P_{n,m}$ is the power for the \textit{m}th subpath of the \textit{n}th path, $\theta_{n,m,\mu}$ is defined as

\begin{equation}
\theta_{n,m,\mu} = \mathrm{mod}\!\left(\theta_{n,m} - \mu_{\theta} + \pi,\; 2\pi\right) - \pi\,,
\label{eq:A2}
\end{equation}

$\mu_{\theta}$ is defined as

\begin{equation}
\mu_{\theta} = \frac{\displaystyle\sum_{n=1}^{N}\sum_{m=1}^{M}\theta_{n,m}\cdot P_{n,m}}{\displaystyle\sum_{n=1}^{N}\sum_{m=1}^{M}P_{n,m}}
\label{eq:A3}
\end{equation}

and $\theta_{n,m}$ is the AoA (or AoD) of the \textit{m}th subpath of the \textit{n}th path. Note that we have dropped the AoA (AoD) subscript for convenience.

We note that the angle spread should be independent of a linear shift in the AoAs. In other words, by replacing $\theta_{n,m}$ with $\theta_{n,m}+\Delta$, the angle spread $\sigma_{AS}(\Delta)$ which is now a function of $\Delta$ should actually be constant no matter what $\Delta$ is. However, due to the ambiguity of the modulo $2\pi$ operation, this may not be the case. Therefore the angle spread should be the minimum of $\sigma_{AS}(\Delta)$ over all $\Delta$:

\begin{equation}
\sigma_{AS} = \min_{\Delta}\,\sigma_{AS}(\Delta) = \sqrt{\frac{\displaystyle\sum_{n=1}^{N}\sum_{m=1}^{M}\bigl(\theta_{n,m,\mu}(\Delta)\bigr)^{2}\cdot P_{n,m}}{\displaystyle\sum_{n=1}^{N}\sum_{m=1}^{M}P_{n,m}}}
\label{eq:A4}
\end{equation}

where $\theta_{n,m,\mu}(\Delta)$ is defined as

\begin{equation}
\theta_{n,m,\mu}(\Delta) = \mathrm{mod}\!\left(\theta_{n,m}(\Delta) - \mu_{\theta}(\Delta) + \pi,\; 2\pi\right) - \pi\,,
\label{eq:A5}
\end{equation}

$\mu_{\theta}(\Delta)$ is defined as

\begin{equation}
\mu_{\theta}(\Delta) = \frac{\displaystyle\sum_{n=1}^{N}\sum_{m=1}^{M}\theta_{n,m}(\Delta)\cdot P_{n,m}}{\displaystyle\sum_{n=1}^{N}\sum_{m=1}^{M}P_{n,m}}
\label{eq:A6}
\end{equation}

and $\theta_{n,m}(\Delta) = \mathrm{mod}\!\left(\theta_{n,m} + \Delta + \pi,\; 2\pi\right) - \pi$.

\end{document}
